\chapter{大数}\label{ch:g4:numbers}

在\cref{ch:g3:numbers}中，数的上限是千；现在，
上限提高到\emph{百万}。读大数的诀窍
是把数字从右往左每三位分成一组，即“级”。

\section{每三位分一级}

\begin{definition}[级]\label{def:g4:numbers:classes}
从右往左数，一个数的数字每三位分为一组，
即\emph{级}：\emph{个}级、
\emph{千}级、\emph{百万}级。每一级都有各自的
个位、十位和百位。例如
\[
6\,309\,452 = 6 \text{ 百万} + 309 \text{ 千}
+ 452 \text{ 个},
\]
读作“六百三十万九千四百
五十二”。
\end{definition}

\begin{omfigure}
\begin{tikzpicture}[scale=0.86]
  \draw (0,0) grid[xstep=1.15, ystep=0.75] (10.35,1.5);
  \node[font=\footnotesize, omThm] at (1.725,1.85) {百万};
  \node[font=\footnotesize, omDef] at (5.175,1.85) {千};
  \node[font=\footnotesize, omMeth] at (8.625,1.85) {个};
  \foreach \x/\l in {0.575/h, 1.725/t, 2.875/u,
                     4.025/h, 5.175/t, 6.325/u,
                     7.475/h, 8.625/t, 9.775/u}
    \node[font=\footnotesize] at (\x,1.12) {\l};
  \foreach \x/\d in {2.875/6, 4.025/3, 5.175/0, 6.325/9,
                     7.475/4, 8.625/5, 9.775/2}
    \node[font=\large] at (\x,0.37) {$\d$};
  \draw[omThm, thick] (0,0) rectangle (3.45,1.5);
  \draw[omDef, thick] (3.45,0) rectangle (6.9,1.5);
  \draw[omMeth, thick] (6.9,0) rectangle (10.35,1.5);
\end{tikzpicture}

{\small $6\,309\,452$ 的数级表（h $=$ 百位，t $=$
十位，u $=$ 各级的个位）。读大数，就是逐级读出，
每级之后读出其名称。}
\end{omfigure}

\begin{example}[用级写数]\label{ex:g4:numbers:write}
读作“一千二百零四万零五百”的数，百万级为 $12$，
千级为 $040$，个级为 $500$，
写作：
\[
12\,040\,500 .
\]
$040$ 中的零不能省略：除最左边一级外，每一级
都必须写足三位数字。
\end{example}

\begin{example}[有几个千？]\label{ex:g4:numbers:contained}
数 $6\,309\,452$ 中\emph{包含} $6\,309$ 个千
（即个级左边的全部数字）：$6\,309\,452 = 6\,309 \times
1\,000 + 452$。问“有几个千”，问的不是
一位数字，而是数的一整段前缀。
\end{example}

\section{比较与四舍五入}

\begin{method}[比较大数]\label{met:g4:numbers:compare}
与\cref{met:g3:numbers:compare}完全相同，只是借助各级
作为捷径：位数多的数较大；位数相同时，从左往右比较——
先逐级比较，级内再逐位比较。
\end{method}

\begin{example}\label{ex:g4:numbers:compare}
$2\,145\,890$ 与 $2\,148\,540$：百万级相同（均为 $2$），千级
$145 < 148$，所以
\[
2\,145\,890 < 2\,148\,540 .
\]
\end{example}

\begin{definition}[四舍五入]\label{def:g4:numbers:round}
将一个数\emph{四舍五入}\index{四舍五入}到最接近的千
（或十、百、百万\,\dots），即换成离它最近的
$1\,000$ 的倍数。恰好位于中间时，进一。实际
操作中：看被取位的下一位数字——$5$ 或
更大则进一，$4$ 或更小则舍去。
\end{definition}

\begin{example}\label{ex:g4:numbers:round}
$37\,614$ 取到最接近的千：百位数字是
$6$，故进一，得 $38\,000$。取到最接近的百：
得 $37\,600$。取到最接近的万：$40\,000$。
\end{example}

\begin{omfigure}
\begin{tikzpicture}[scale=1.05]
  \draw[-Stealth] (-0.3,0) -- (10.8,0);
  \draw (0,-0.1) -- (0,0.1);
  \node[below, font=\small] at (0,-0.12) {$37\,000$};
  \draw (10,-0.1) -- (10,0.1);
  \node[below, font=\small] at (10,-0.12) {$38\,000$};
  \draw (5,-0.08) -- (5,0.08);
  \node[below, font=\footnotesize] at (5,-0.12) {$37\,500$};
  \fill[omThm] (6.14,0) circle (2pt)
    node[above, font=\small] {$37\,614$};
  \draw[-Stealth, omDef, thick] (6.14,0.42) .. controls (8,0.85) ..
    (9.9,0.28);
\end{tikzpicture}

{\small 将 $37\,614$ 取到最接近的千：它已超过
中点 $37\,500$，故进一，得 $38\,000$。}
\end{omfigure}

\begin{example}[世界中的数量级]\label{ex:g4:numbers:world}
\omterm{def:g4:numbers:round}{四舍五入}是大数在实际中的用途：某国有
$67\,842\,591$ 人，可说成“约 $68$ 百万”；某城市有
$214\,376$ 人，可说成“约 $200\,000$”。精确数字本来每天都在
变化，\omterm{def:g4:numbers:round}{四舍五入}后的数字
才携带有用的信息。
\end{example}

\section{练习}

\begin{exercise}[$\star$]\label{exo:g4:numbers:1}
用数字写出（数级之间用逗号分开）：
“三百五十二万”；“七百万零八”；
“二十四万零六十”。
\end{exercise}

\begin{exercise}[$\star$]\label{exo:g4:numbers:2}
用文字写出：$5\,204\,300$；\quad $10\,050\,000$；\quad
$999\,999$。
\end{exercise}

\begin{exercise}[$\star$]\label{exo:g4:numbers:3}
在 $8\,472\,156$ 中：万位上是哪个数字？百万级
中的百位上是哪个数字？这个数含有多少个\emph{千}
（见\cref{ex:g4:numbers:contained}）？
\end{exercise}

\begin{exercise}[$\star$]\label{exo:g4:numbers:4}
抄写并在空处填上 $<$ 或 $>$：
\[
980\,000 \;?\; 1\,020\,000, \qquad
3\,456\,000 \;?\; 3\,465\,000, \qquad
7\,090\,050 \;?\; 7\,090\,500 .
\]
\end{exercise}

\begin{exercise}[$\star$]\label{exo:g4:numbers:5}
按从小到大的顺序排列：$2\,300\,000$；\ $980\,500$；\
$2\,030\,000$；\ $2\,299\,999$。
\end{exercise}

\begin{exercise}[$\star$]\label{exo:g4:numbers:6}
对 $63\,782$ 作\omterm{def:g4:numbers:round}{四舍五入}，分别取到最接近的十、百、
千、万。
\end{exercise}

\begin{exercise}[$\star$]\label{exo:g4:numbers:7}
\omterm{def:g4:numbers:round}{四舍五入}到最接近的百万：$4\,721\,000$；\quad $6\,499\,999$；
\quad $9\,500\,000$。
\end{exercise}

\begin{exercise}[$\star$]\label{exo:g4:numbers:8}
按步计数：以 $1\,000$ 为步长，从 $998\,500$ 出发数三步；
以 $100\,000$ 为步长，从 $950\,000$ 出发数两步。越过
一百万时，各数级会发生什么变化？
\end{exercise}

\begin{exercise}[$\star$]\label{exo:g4:numbers:9}
某体育场可容纳 $81\,338$ 名观众。某报纸写“超过
$80\,000$ 名球迷”，某电台说“约 $81\,000$”。两种说法
都对吗？各是\omterm{def:g4:numbers:round}{四舍五入}到哪一位？
\end{exercise}

\begin{exercise}[$\star\star$]\label{exo:g4:numbers:10}
用数字 $1$ 至 $7$ 每个恰好使用一次，写出
最大的七位数、最小的七位数，以及最接近
$4\,000\,000$ 的那个数。
\end{exercise}

\begin{exercise}[$\star\star$]\label{exo:g4:numbers:11}
若\omterm{def:g4:numbers:round}{四舍五入}到最接近的千，结果为 $50\,000$ 的
最大整数是多少？最小的呢？
\end{exercise}
